% ===========================================================================
\section{Suprema and chaining}\label{sec:chaining}
\warmth{0}\quad\whisper{The exit. When the target is $\sup_t X_t$ (an empirical process), single-point Chernoff + a union bound is too loose.}

\begin{strand}
In one line: the first seven sections control \emph{one} random quantity. But ML/statistics really wants
$\sup_{t\in T}X_t$ (empirical processes, the sup in a generalization bound). \keyword{Chaining} uses multi-scale nets
to control the supremum by an integral of metric entropy, reducing ``infinitely many tails'' to ``the geometry of
covering numbers''. This is the bridge to empirical processes.
\end{strand}

\block{The paradigm}
\trigger{You want $\E\sup_{t\in T}X_t$, where $\{X_t\}$ is a (sub-)Gaussian process with natural metric $d(s,t)=\|X_s-X_t\|_{\psi_2}$.}
\keyword{Chaining}: along a sequence of refining $\varepsilon$-nets $T_0\subset T_1\subset\cdots$, telescope $X_t-X_{t_0}$ and sum scale by scale. \keyword{Dudley}:
\[
  \E\sup_{t\in T}X_t\ \le\ C\int_0^\infty\sqrt{\log N(T,d,\varepsilon)}\ d\varepsilon.
\]

\begin{strand}
Three-line heuristic: (1) a union bound $=$ a single scale, chaining $=$ \emph{all scales} summed, hence tighter;
(2) the currency is \keyword{metric entropy} $\log N(T,d,\varepsilon)$, turning a probability problem into covering geometry;
(3) Dudley can be loose; the \keyword{generic chaining} quantity $\gamma_2(T,d)$ is the \emph{tight} one (Talagrand's majorizing measures).
\end{strand}

\block{Cheat-sheet: from union bound to majorizing measures}
\noindent\begin{tabularx}{\linewidth}{@{}l l L@{}}
\toprule
{\color{indigo}Tool} & {\color{indigo}Gives} & \multicolumn{1}{l}{\color{indigo}When to use / tight?}\\
\midrule
Union bound     & $\sqrt{2\log|T|}$ (finite $T$)  & opening move; too loose for large or continuous $T$\\
Dudley integral & $\int\sqrt{\log N(\varepsilon)}\,d\varepsilon$ & standard upper bound for continuous index; often enough\\
Sudakov lower bd.\ & $\sup_\varepsilon\varepsilon\sqrt{\log N(\varepsilon)}$ & pairs with Dudley to pin the order\\
Generic chaining& $\gamma_2(T,d)$               & matching upper \emph{and} lower (Gaussian processes, Talagrand)\\
\bottomrule
\end{tabularx}

\block{Main results \& proof skeletons}

\begin{theorem}[Dudley's inequality]\label{thm:dudley}
$\{X_t\}_{t\in T}$ a centered process, sub-Gaussian w.r.t.\ $d$. Then
$\E\sup_{t}X_t\le C\int_0^{\infty}\sqrt{\log N(T,d,\varepsilon)}\,d\varepsilon$.
\end{theorem}
\begin{proof}
Skeleton: take nets $T_k$ at scale $\varepsilon_k=2^{-k}$ with nearest-point maps $\pi_k(t)$; telescope
$X_t-X_{\pi_0(t)}=\sum_{k\ge1}(X_{\pi_k(t)}-X_{\pi_{k-1}(t)})$; each link has $|T_k||T_{k-1}|$ sub-Gaussian increments,
combine the per-scale union bound with the increment scale $\lesssim 2^{-k}$ and sum.
\TODO{fill the per-scale union bound $\sqrt{\log|T_k|}$ and the geometric-series sum $\Rightarrow$ the integral}
\end{proof}

\begin{theorem}[Generic chaining / majorizing measures]\label{thm:generic-chaining}
For a Gaussian process, $\E\sup_t X_t\asymp\gamma_2(T,d)$, where
$\gamma_2(T,d)=\inf_{(T_k)}\sup_t\sum_{k\ge0}2^{k/2}d(t,T_k)$ with $|T_k|\le 2^{2^k}$.
\end{theorem}
\begin{proof}\TODO{upper bound: weighted chaining (weights adapt to $t$, beating Dudley's uniform scales); lower bound: Talagrand's partition tree / iterated Sudakov. Cite the majorizing-measures theorem (do not reprove)}\end{proof}

\begin{example}[Operator norm via $\varepsilon$-net + chaining]\label{ex:operator-norm}
$\|A\|=\sup_{u,v\in S^{n-1}}u^\top Av$: discretize on a $1/2$-net of the sphere ($\le 9^n$ points) + each point sub-Gaussian
(§\ref{sec:subgaussian}) $\Rightarrow$ $\E\|A\|\lesssim\sqrt n$. Chaining ``automates'' the net argument to a continuous index.
\TODO{fill the standard $1/2$-net lemma lifting $\sup$ to $2\max_{\text{net}}$}
\end{example}

\begin{strand}
Why it goes this way: the Borell–TIS of §\ref{sec:geometry} (\pick{borelltis}) controls the \emph{fluctuation} of $\sup_t X_t$,
this section its \emph{expectation}, together they are the full concentration of an empirical process. This closes the loop
with the ``expectation $=$ Rademacher complexity'' half of the generalization bound in §\ref{sec:bdd}.
\end{strand}

\loose{This is the exit: metric entropy / empirical processes / uniform laws deserve \emph{their own notebook} (Dudley, van der Vaart–Wellner, Talagrand's \emph{Upper and Lower Bounds}). This section only goes as far as ``chaining + Borell–TIS'', as a bridge.}%
\warp{chaining} Chaining $\to$ Rademacher complexity $\to$ uniform laws of large numbers.
